\documentclass[12pt]{article}
\usepackage[utf8]{inputenc}
\usepackage[T1]{fontenc}
\usepackage[margin=1in]{geometry}
\usepackage{setspace}
\usepackage{xcolor}
\usepackage[colorlinks=true, urlcolor=blue!60!black, linkcolor=blue!60!black]{hyperref}
\setlength{\parindent}{2em}
\setlength{\parskip}{1.25em}
\renewcommand{\baselinestretch}{1.0}
\begin{document}
\begin{center}
{\LARGE \textit{Citizens United v. Federal Election Commission}, 558 U.S. 310 (2010)\par}
\vspace{0.8em}
{\large Case Brief --- Four-Page Edition\par}
\end{center}
\vspace{1em}

\section{Memory Jogger}

A 5--4 Court struck down federal bans on corporate and union independent political expenditures under the First Amendment, overruling \textit{Austin}, while upholding disclaimer and disclosure requirements 8--1 (\href{https://www.courtlistener.com/opinion/1741/citizens-united-v-federal-election-commission/}{CourtListener}).

\section{Detailed Case Facts}

Citizens United, a conservative nonprofit corporation funded mostly by individual donors and in small part by for-profit corporations, released \textit{Hillary: The Movie} in January 2008, a 90-minute documentary sharply criticizing then-Senator Hillary Clinton during the Democratic presidential primaries. It sought to pay a cable provider \$1.2 million to offer the film free on video-on-demand to 34.5 million subscribers and to air short television ads promoting it.

Section 441b, as amended by \S~203 of the Bipartisan Campaign Reform Act of 2002 (BCRA, or McCain--Feingold), barred corporations and unions from using general treasury funds for express advocacy or any ``electioneering communication''---a broadcast, cable, or satellite communication naming a federal candidate within 30 days of a primary or 60 days of a general election. Corporations could still speak through a ``separate segregated fund'' (a PAC). BCRA \S\S~201 and 311 required on-air disclaimers and FEC disclosure. \textit{Buckley v. Valeo} (1976) had struck down independent-expenditure ceilings, but \textit{Austin v. Michigan Chamber of Commerce} (1990) upheld a corporate-expenditure ban on an ``antidistortion'' rationale, and \textit{McConnell} (2003) relied on \textit{Austin} to uphold \S~203.

\section{Procedural History}

Citizens United sued the FEC in December 2007, challenging \S~441b as applied to \textit{Hillary} and the disclosure provisions as applied to the film and its ads; it stipulated to dismissing its facial challenge. A three-judge district court denied a preliminary injunction and then granted the FEC summary judgment, finding the film the ``functional equivalent'' of express advocacy. On direct appeal, after argument in March 2009, the Supreme Court ordered reargument on whether to overrule \textit{Austin} and part of \textit{McConnell}, and decided the case on January 21, 2010.

\section{Judicial Votes}

\begin{itemize}
\item \textbf{Majority:} Justice Kennedy, joined by Roberts, Scalia, Alito, and (except Part IV) Thomas: 5--4 on the expenditure ban. Stevens, Ginsburg, Breyer, and Sotomayor joined Part IV, making it 8--1 on disclosure.
\item \textbf{Concurrences:} Chief Justice Roberts (with Alito); Justice Scalia (with Alito, and Thomas in part).
\item \textbf{Concurring in part and dissenting in part:} Justice Stevens (with Ginsburg, Breyer, and Sotomayor); Justice Thomas (dissenting from Part IV).
\end{itemize}

\section{Holding}

The government may not suppress political speech on the basis of the speaker's corporate identity. Bans on corporate (and union) independent expenditures fail strict scrutiny because such expenditures do not give rise to \textit{quid pro quo} corruption or its appearance; \textit{Austin} and \textit{McConnell}'s \S~203 holding are overruled. Disclaimer and disclosure requirements survive ``exacting scrutiny'' on the informational interest, absent a reasonable probability of threats or reprisals. The Court left open restrictions on foreign nationals and did not disturb limits on direct contributions.

\section{Analysis of Opinions}

\textbf{Majority (Kennedy).} The Court rejected narrower statutory and as-applied routes, reasoning that a complex FEC regime operated as a \textit{de facto} prior restraint. Section 441b was an outright ban backed by felony penalties; the PAC option did not cure it. The antidistortion rationale fails because government has no interest in equalizing voices, and its logic would let Congress censor media corporations. Corruption means \textit{quid pro quo} exchanges; ingratiation and access are not corruption. Shareholder protection was both underinclusive and overinclusive, and dissenters can rely on ``the procedures of corporate democracy.'' \textit{Stare decisis} did not save \textit{Austin}. Part IV upheld disclosure as imposing no ceiling on speech.

\textbf{Roberts (concurring).} ``There is a difference between judicial restraint and judicial abdication''; \textit{Austin} was an aberration that could not be saved by ``jury-rigging new and different justifications.''

\textbf{Scalia (concurring).} The First Amendment protects speech, not categories of speakers, and founding-era hostility targeted state-granted monopolies, not modern corporations.

\textbf{Stevens (dissenting in part).} The majority revived an abandoned facial challenge without a record. Section 441b was not a ``ban,'' since corporations could speak through PACs or outside the windows. Corporations cannot vote or run for office, and the Framers understood corruption far more broadly than bribery.

\textbf{Thomas (dissenting in part).} Disclosure exposed Proposition 8 donors to threats, boycotts, and reprisals; slow as-applied challenges cannot protect them.

\section{Examples: Future Applications}

\begin{enumerate}
\item \textbf{Same-Side (Invalidated): State Corporate-Expenditure Ban.} A state bar on corporate treasury spending on independent ads in legislative races is unconstitutional, as \textit{American Tradition Partnership v. Bullock} (2012) summarily held.
\item \textbf{Same-Side (Invalidated): Super PAC Contribution Caps.} Caps on contributions to committees that make only independent expenditures lack a corrupting nexus, as \textit{SpeechNow.org v. FEC} (D.C. Cir. 2010) held.
\item \textbf{Opposite-Side (Upheld): Direct Corporate Contributions.} The ban on corporate contributions to candidates survives, because direct contributions pose a concrete risk of \textit{quid pro quo} corruption.
\item \textbf{Opposite-Side (Upheld): Digital Ad Disclaimers.} A state rule requiring sponsor disclaimers and donor reports for corporate online election ads survives exacting scrutiny absent proof of likely harassment.
\item \textbf{Fence-Sitter: Foreign-Owned Subsidiaries.} A ban on independent spending by U.S. subsidiaries majority-owned by foreigners pits the principle against speaker discrimination against the interest in excluding foreign nationals from self-government, which the Court left open.
\end{enumerate}

\section{Critique}

\textbf{Progressive and originalist critics.} Richard Hasen and the Brennan Center argue that the claim that independent expenditures do not corrupt is contradicted by reality: single-candidate Super PACs blur the line with coordinated spending, and ``dark money'' groups obscure donors. Zephyr Teachout argues that the Framers' concept of corruption reached improper dependence, and other scholars stress that founding-era corporations were narrowly chartered creatures of state law.

\textbf{Free-speech defenders.} Joel Gora praises the decision for protecting millions of small businesses and nonprofits, from the ACLU to the NRA, and notes the government's concession that \S~441b could theoretically ban books. At bottom, the case pits a libertarian view of the First Amendment against a republican view that legislatures may protect electoral equality from concentrated wealth.

\section{Key Quotations}

\begin{quote}
``For the reasons explained above, we now conclude that independent expenditures, including those made by corporations, do not give rise to corruption or the appearance of corruption.''
\end{quote}
\noindent --- Justice Kennedy, Opinion of the Court, 558 U.S. at 357.

\begin{quote}
``We return to the principle established in \textit{Buckley} and \textit{Bellotti} that the Government may not suppress political speech on the basis of the speaker's corporate identity. No sufficient governmental interest justifies limits on the political speech of nonprofit or for-profit corporations.''
\end{quote}
\noindent --- Justice Kennedy, Opinion of the Court, 558 U.S. at 365.

\begin{quote}
``The Government may regulate corporate political speech through disclaimer and disclosure requirements, but it may not suppress that speech altogether.''
\end{quote}
\noindent --- Justice Kennedy, Opinion of the Court, 558 U.S. at 319.

\end{document}
